\section{Conclusão}
\label{sec:conclusao}

Construímos um removedor de fundo em tempo real para webcam e vídeos,
sem redes neurais, combinando a subtração de fundo com MOG2 e as
técnicas vistas em aula. O MOG2 decide, pixel a pixel, o que é pessoa,
e as outras etapas limpam e completam essa decisão. O sistema acertou a
contagem de objetos em 96,2\% dos quadros e chegou a um IoU médio de
0,791, com 4 de 12 quadros acima do critério de 0,8.

A principal lição foi que quase todos os erros vieram de quebrar um
pressuposto do método, e não do algoritmo: a câmera que muda a
exposição sozinha, a pessoa na cena durante o aprendizado e a roupa da
cor do fundo. O exemplo feito à mão mostrou por que a exposição era tão
grave e por que aumentar o limiar não resolvia. Aprendemos também a
medir antes de mudar: medir o brilho da parede revelou a causa, e a
melhor correção para a webcam foi na captura, travando a exposição.
Por fim, vimos que todo ajuste é uma troca e que até o gabarito precisa
ser conferido antes de se confiar no número.
